% !TeX program = xelatex
\documentclass[UTF8,fontset=fandol,11pt,a4paper]{ctexart}
\usepackage[margin=23mm]{geometry}
\usepackage{amsmath,amssymb,amsthm,mathtools,bm}
\usepackage{booktabs,longtable,tabularx,array}
\usepackage{xcolor,graphicx}
\usepackage[colorlinks=true,linkcolor=blue!50!black,urlcolor=blue!60!black,citecolor=blue!50!black]{hyperref}
\usepackage{enumitem}
\setlist{itemsep=3pt,topsep=4pt}
\allowdisplaybreaks[3]
\newtheorem{proposition}{命题}[section]
\newtheorem{lemma}[proposition]{引理}
\newtheorem{theorem}[proposition]{定理}
\theoremstyle{remark}
\newtheorem{remark}[proposition]{说明}
\newcommand{\T}{\mathbb T}
\newcommand{\R}{\mathbb R}
\newcommand{\eps}{\varepsilon}
\newcommand{\dd}{\,\mathrm d}
\newcommand{\avg}[1]{\overline{#1}}
\newcommand{\norm}[1]{\left\|#1\right\|}
\DeclareMathOperator{\diver}{div}
\title{仅水平黏性的二维等熵可压缩流体\\二阶声学生成、证明路线与文献核对}
\author{研究备忘录：第一轮独立评估}
\date{2026年10月30日}
\begin{document}
\maketitle
\begin{abstract}
本文核对固定退化黏性模型的一维激波嵌入、二维涡旋初值产生声学零模的二阶机制，并区分精确恒等式、可严格化的有限时间展开与长时间激波猜想。推荐先证明可检验的有效声学动力学及误差定理。
\end{abstract}
\tableofcontents
\section{研究判断与结论层次}

原始的存在性问题已有一个直接的肯定答案：系统含有一维等熵 Euler 的不变子类，压缩简单波给出有限时间梯度爆破。这一事实本身不宜作为论文的主要创新点。更有内容的问题是：\emph{初始没有声学零模、只有真正二维涡旋振荡时，耗散过程中产生的二次声学波如何被选择，如何在后续时间中演化？}

本轮计算支持继续研究，但支持的范围必须限定如下。
\begin{enumerate}
\item \textbf{精确事实：}一维嵌入、瞬时零模生成、精确平均方程及线性谱都可以直接验证。
\item \textbf{展开系数的精确事实：}给定测试初值的一阶密度为零；二阶水平平均包含非零、无阻尼的自由声学部分，对任何固定 $\mu,A>0$、$\gamma>1$ 均无完全抵消。
\item \textbf{能量估计结论：}本文给出固定时间与对数时间二阶近似的定量误差证明方案，并由小数据能量重启得到寿命下界；作为投稿稿件，仍需完整整理局部构造和各阶稳定性引理。
\item \textbf{尚未证明：}完整二维解在 $O(\eps^{-2})$ 时间上的有效动力学、激波形成及寿命渐近。二阶系数在 $t\to\infty$ 时的行为并不等于固定 $\eps$ 的原方程长期行为。
\end{enumerate}
全文固定 $A,\gamma,\mu$，令 $\eps\to0$。所有误差常数允许依赖这三个参数，不声称对 $\mu\to0$、$\mu\to\infty$ 或声速趋零一致。

\section{模型、能量与一维激波嵌入}
记 $x=x_1,y=x_2$，$c=\sqrt{p'(1)}=\sqrt{A\gamma}$。研究且仅研究
\begin{equation}\label{eq:model}
 \rho_t+\diver(\rho u)=0,\qquad
 (\rho u)_t+\diver(\rho u\otimes u)+\nabla p(\rho)=\mu\partial_x^2u,
 \quad (x,y)\in\T^2.
\end{equation}
其中，$t\ge0$ 是时间，$x=x_1$ 是水平方向坐标，$y=x_2$ 是垂直方向坐标；$\rho=\rho(t,x,y)>0$ 是密度，$u=(u_1,u_2)$ 是速度，$u_1,u_2$ 分别是水平、垂直速度分量；$p(\rho)=A\rho^\gamma$ 是等熵压力律，$A>0$ 是压力常数，$\gamma>1$ 是绝热指数，$\mu>0$ 是仅作用于水平方向的动力黏性系数。$\T=\R/(2\pi\mathbb Z)$ 是长度为 $2\pi$ 的周期区间，$\T^2=\T\times\T$ 是二维环面，$\R$ 和 $\mathbb Z$ 分别是实数集和整数集。

其中，$\nabla=(\partial_x,\partial_y)$ 是空间梯度，$\diver$ 是空间散度，$u\otimes u$ 是分量为 $u_i u_j$ 的张量积，张量的散度按动量分量计算；$\partial_x^2u$ 表示对两个速度分量分别取水平二阶导数。下标 $t,x,y$ 表示对相应变量求偏导，初值下标 $0$ 表示 $t=0$ 时的取值。
正密度的光滑解等价地满足
\begin{equation}\label{eq:primitive}
 u_t+u\cdot\nabla u+\frac{p'(\rho)}{\rho}\nabla\rho
 =\frac{\mu}{\rho}\partial_x^2u.
\end{equation}
其中，$u\cdot\nabla=u_1\partial_x+u_2\partial_y$ 是沿速度方向的输运算子，$p'(\rho)$ 是压力对密度的导数；$c=\sqrt{p'(1)}$ 是常密度平衡态 $\rho=1$ 处的声速，$\mu/\rho$ 是速度方程中的可变扩散系数。
所以不能把速度方程中的扩散系数当作常数，也不能附加 $\diver u=0$。
周期边界给出质量和两个方向总动量守恒，以及
\begin{equation}\label{eq:energy}
 \frac{\mathrm d}{\mathrm dt}\int_{\T^2}
 \left(\frac12\rho|u|^2+\frac{A}{\gamma-1}\rho^\gamma\right)\dd x\dd y
 +\mu\int_{\T^2}|\partial_xu|^2\dd x\dd y=0.
\end{equation}
其中，$|u|$ 是速度向量的欧氏长度，$\dd x\dd y$ 是环面上的面积元；$\rho|u|^2/2$ 是动能密度，$A\rho^\gamma/(\gamma-1)$ 是内能密度，$\mu\int_{\T^2}|\partial_xu|^2\dd x\dd y$ 是水平黏性造成的能量耗散率。
这不是对 $\partial_yu$ 的耗散估计；能量有界与梯度爆破相容。

\subsection{约化完全成立，但只在不变子类内}
代入 $\rho=r(t,y)$、$u=(0,v(t,y))$ 后，水平黏性和水平动量方程消失，余下
\begin{equation}\label{eq:1d}
 r_t+(rv)_y=0,\qquad (rv)_t+(rv^2+p(r))_y=0.
\end{equation}
其中，$r=r(t,y)$ 是该一维不变子类中的密度，$v=v(t,y)$ 是垂直速度，$rv$ 是垂直动量密度；$r_0=r(0,\cdot)$ 和 $v_0=v(0,\cdot)$ 分别是它们的初值。
需要 $r_0>0$、$r_0,v_0$ 光滑且 $2\pi$ 周期。由光滑解唯一性，该嵌入在首次奇性前成立；不要求初始密度常数。

令 $c(r)=\sqrt{A\gamma}\,r^{(\gamma-1)/2}$，黎曼不变量为
\[
 w_\pm=v\pm\frac{2c(r)}{\gamma-1},\qquad
 (\partial_t+(v\pm c(r))\partial_y)w_\pm=0.
\]
其中，$c(r)=\sqrt{p'(r)}$ 是局部密度为 $r$ 时的声速，$w_+,w_-$ 是两支声学特征族的黎曼不变量，$v+c(r)$ 和 $v-c(r)$ 是对应的特征速度；符号 $\pm$ 在同一式中按相同的正号或负号配对使用。
固定 $w_-=K$，则
\begin{equation}\label{eq:simple}
 (w_+)_t+\left(\frac{\gamma+1}{4}w_+
       +\frac{3-\gamma}{4}K\right)(w_+)_y=0.
\end{equation}
其中，$K$ 是在整个时空中固定的 $w_-$ 值；只让另一不变量 $w_+$ 变化便得到这里的简单波。
例如取 $K=-2c/(\gamma-1)$、
$w_+(0,y)=2c/(\gamma-1)+\delta\sin y$，$\delta>0$ 充分小。于是
\[
 v_0=\frac\delta2\sin y,\qquad
 c(r_0)=c+\frac{\gamma-1}{4}\delta\sin y>0.
\]
其中，$\delta>0$ 是这个一维示例的波幅参数，用其充分小性保证初始声速和密度为正；它与后文二维测试初值的振幅参数 $\eps$ 分别使用。
特征映射及斜率公式为
\[
 y=\xi+\lambda_+(w_{+,0}(\xi),K)t,\qquad
 (w_+)_y(t,y)=\frac{w_{+,0}'(\xi)}{1+\frac{\gamma+1}{4}t w_{+,0}'(\xi)}.
\]
其中，$\xi\in\T$ 是特征曲线在 $t=0$ 时的出发位置，$w_{+,0}(\xi)=w_+(0,\xi)$ 是初始黎曼不变量，$w_{+,0}'$ 是它对初始位置的导数；$\lambda_+(w,K)=(\gamma+1)w/4+(3-\gamma)K/4$ 是该简单波的正族特征速度。
故精确首次激波时间
\begin{equation}\label{eq:embeddedtime}
 T_*=-\frac4{(\gamma+1)\min w_{+,0}'}=\frac4{(\gamma+1)\delta}.
\end{equation}
其中，$T_*$ 是这个嵌入式一维解首次出现无限斜率的时间，$\min w_{+,0}'$ 是初始斜率在整个周期区间上的最小值；它在此正弦示例中等于 $-\delta$。
这里 $v,r$ 本身一直有界，而且
\[
 0<r_{\min}\le r(t,y)\le r_{\max}<\infty.
\]
其中，$r_{\min},r_{\max}$ 可分别取初始密度的最小值和最大值。这一界成立，是因为 $r,v$ 是沿特征搬运的 $w_+$ 的光滑函数。爆破的是 $v_y$、$r_y$ 的 $L^\infty$ 范数；这里 $L^\infty$ 范数是函数绝对值的本质上确界。密度梯度与 $w_+$ 梯度间的系数在上述紧正区间内非零。原二维解的 $C^1$ 正则性因此失效，$C^1$ 表示函数具有连续的一阶导数。这里不涉及真空或速度幅值无穷大，也未证明任意真正二维扰动下该激波稳定。

\subsection{初始密度常数不构成不变约束}
对于本题测试初值，$\rho_t(0)=0$，但直接微分得到
\begin{equation}\label{eq:rhott}
 \rho_{tt}(0,x,y)=\eps^2(\cos2x-\cos2y).
\end{equation}
其中，$\rho_{tt}$ 是密度对时间的二阶偏导数，$\eps>0$ 是二维测试初值的速度振幅参数；式中取值时刻为 $t=0$，并使用了 $\rho_0=1$。
因此密度立即在二阶时间展开中偏离常数。初始散度零只用于计算初始及一阶量。

\section{线性谱与测试初值的一阶解}
采用不带阶乘的展开约定
\begin{equation}\label{eq:series}
 \rho=1+\eps a+\eps^2b+\cdots,\qquad
 u=\eps U+\eps^2 Z+\cdots.
\end{equation}
其中，$a,b$ 分别是密度关于 $\eps$ 的一阶、二阶展开系数，$U=(U_1,U_2)$ 和 $Z=(Z_1,Z_2)$ 分别是速度的一阶、二阶展开系数，均为时空函数；省略号表示更高阶项。系数中未吸收 $\eps$ 的幂次，二阶项也未另乘 $1/2$。
一阶方程为
\begin{equation}\label{eq:linear}
 a_t+\diver U=0,\qquad U_t+c^2\nabla a=\mu\partial_x^2U.
\end{equation}
频率 $k=(k_1,k_2)\ne0$ 处，横于 $k$ 的速度分量有特征值 $-\mu k_1^2$；纵向声学特征值满足
\begin{equation}\label{eq:spectrum}
 \lambda^2+\mu k_1^2\lambda+c^2|k|^2=0,
 \quad \lambda_\pm=\frac{-\mu k_1^2\pm\sqrt{\mu^2 k_1^4-4c^2|k|^2}}2.
\end{equation}
其中，$k=(k_1,k_2)\in\mathbb Z^2$ 是空间 Fourier 频率，$k_1,k_2$ 分别是水平、垂直波数，$|k|=\sqrt{k_1^2+k_2^2}$；$\lambda$ 是模式 $e^{\lambda t+i(k_1x+k_2y)}$ 的时间增长率，$\lambda_\pm$ 是声学部分的两个特征根，$i$ 是满足 $i^2=-1$ 的虚数单位。这里的 $\lambda_\pm$ 是线性谱值，区别于前面一维简单波的特征速度 $\lambda_+(w,K)$。
当 $k_1=0,k_2\ne0$，有 $\lambda_\pm=\pm ic|k_2|$，以及水平剪切特征值 $0$。因此“水平零模”包含两支声学模和一个剪切模，不能全部称为声学模。$k=(0,0)$ 对应守恒量。
对于 $k_1\ne0$，各根实部为负；固定 $\mu,c$ 的周期格点上可建立指数衰减，但密度衰减需用声学耦合或补偿能量，单凭速度的 Poincar\'e 不等式不够。

测试初值记作
\[
 W=(\sin x\sin y,\ \cos x\cos y),\qquad a(0)=0,\quad U(0)=W.
\]
其中，$W=W(x,y)$ 是不含振幅因子 $\eps$ 的初始速度形状，因此原初值是 $u_0=\eps W$；$a(0)=0$ 表示初始密度没有一阶扰动。
由于 $\diver W=0$ 且 $\partial_x^2W=-W$，唯一一阶解为
\begin{equation}\label{eq:first}
 a\equiv0,\qquad U=e^{-\mu t}W.
\end{equation}
这是一阶特殊结构，不是原方程不可压缩性的证明。

\section{二阶全场与水平平均的显式计算}
\subsection{瞬时生成公式}
先明确水平平均的约定：
\[
 \avg f(t,y)=\frac1{2\pi}\int_0^{2\pi}f(t,x,y)\dd x.
\]
其中，$f$ 是待平均的标量场，向量场按分量取平均；上横线表示只对水平变量 $x$ 积分所得的平均，结果仍依赖 $t,y$，并非对整个二维环面取平均。
平均垂直动量的精确守恒式给出
\[
 \left.\partial_t\avg{\rho u_2}\right|_0
 =-\partial_y\avg{u_{2,0}^2}
 =-\partial_y\left(\frac{\eps^2}2\cos^2y\right)
 =\frac{\eps^2}2\sin2y.
\]
其中，$\avg{\rho u_2}$ 是水平平均的垂直动量密度，$u_{2,0}=\eps\cos x\cos y$ 是初始垂直速度，$\left.\cdot\right|_0$ 表示在 $t=0$ 取值。
用户给出的公式正确。它产生的是 $k_1=0,k_2=\pm2$ 的水平零模，不是总动量 $k=(0,0)$；后者仍为零。

\subsection{二阶方程及被消掉的项}
因为 $a=0$，二阶方程简化为
\begin{equation}\label{eq:second}
 b_t+\diver Z=0,\qquad
 Z_t+c^2\nabla b-\mu\partial_x^2Z=-U\cdot\nabla U,
 \quad (b,Z)|_{t=0}=0.
\end{equation}
一般二阶速度方程中的 $[p''(1)-p'(1)]a\nabla a$、$-\mu a\partial_x^2U$，以及质量方程中的 $\diver(aU)$ 在这里均为零。直接计算
\begin{equation}\label{eq:force}
 -U\cdot\nabla U
 =e^{-2\mu t}\left(-\frac12\sin2x,\frac12\sin2y\right)
 =e^{-2\mu t}\nabla\left(\frac14\cos2x-\frac14\cos2y\right).
\end{equation}
源项虽是梯度，却不能像不可压缩方程的拉格朗日乘子那样被“压力投影”消掉；这里压力由密度决定，恰会激发声学响应。

二阶全场可写成
\begin{equation}\label{eq:fullsecond}
 b=H(t)\cos2x+B(t)\cos2y,\quad
 Z=(P(t)\sin2x,Q(t)\sin2y),
\end{equation}
其中，$H(t),B(t)$ 分别是二阶密度在水平模式 $\cos2x$ 和垂直模式 $\cos2y$ 上的系数，$P(t),Q(t)$ 分别是二阶水平、垂直速度的正弦模式系数；$P(t)$ 不是压力函数 $p(\rho)$。这些系数满足
\begin{align}
 H'=-2P,&\quad P'+4\mu P-2c^2H=-\tfrac12e^{-2\mu t},\label{eq:Hx}\\
 B'=-2Q,&\quad Q'-2c^2B=\tfrac12e^{-2\mu t}.\label{eq:BQ}
\end{align}
其中，$H',B',P',Q'$ 表示各时间系数对 $t$ 的一阶导数，下文 $H'',B''$ 表示其时间二阶导数。
初值均为零。因此
\begin{align}
 H''+4\mu H'+4c^2H&=e^{-2\mu t},\label{eq:HOde}\\
 B''+4c^2B&=-e^{-2\mu t}.\label{eq:BOde}
\end{align}
前者水平受阻尼，后者不受直接阻尼。可选的全场显式公式是
\[
 H(t)=e^{-2\mu t}
 \begin{cases}
 \dfrac{1-\cos(2\sqrt{c^2-\mu^2}\,t)}{4(c^2-\mu^2)},&c>\mu,\\[5pt]
 t^2/2,&c=\mu,\\[3pt]
 \dfrac{\cosh(2\sqrt{\mu^2-c^2}\,t)-1}{4(\mu^2-c^2)},&c<\mu,
 \end{cases}
 \qquad P=-H'/2.
\]
其中，$\cosh$ 是双曲余弦函数；三种分支分别对应 $c>\mu$、$c=\mu$、$c<\mu$ 时水平二阶模式的欠阻尼、临界阻尼和过阻尼情形。
对固定正参数，$H,P$ 指数衰减（临界情形的多项式因子可由稍慢指数吸收）。

\subsection{声学平均公式与非抵消}
设 $D=\mu^2+c^2$。方程 \eqref{eq:BQ} 的解为
\begin{equation}\label{eq:Bexplicit}
 \boxed{B(t)=\frac{\cos(2ct)-\frac\mu c\sin(2ct)-e^{-2\mu t}}{4D}},
\end{equation}
\begin{equation}\label{eq:Qexplicit}
 \boxed{Q(t)=\frac{\mu\cos(2ct)+c\sin(2ct)-\mu e^{-2\mu t}}{4D}}.
\end{equation}
其中，$D=\mu^2+c^2$ 是为简化表达式引入的正常数；$B,Q$ 分别描述水平平均中的二阶密度和二阶垂直速度。下文使用的 $V=\avg{\rho u_2}/\avg\rho$ 是质量加权的平均垂直速度，其精确定义及方程见后文。
在二阶精度下，$\avg\rho=1+\eps^2B\cos2y$、$\avg{\rho u_2}=\eps^2Q\sin2y$，质量加权速度 $V$ 也是 $\eps^2Q\sin2y$。

去掉指数瞬态，记余下 $B_\infty(t),Q_\infty(t)$；下标不表示时间极限，而表示自由声学轨道。令
\[
 z_\pm(t,y)=Q_\infty(t)\sin2y\pm cB_\infty(t)\cos2y.
\]
其中，$B_\infty=B+e^{-2\mu t}/(4D)$、$Q_\infty=Q+\mu e^{-2\mu t}/(4D)$ 是删去指数瞬态后的系数；$z_+,z_-$ 是二阶自由声学态在两支线性特征方向上的组合，尚未乘回 $\eps^2$。
则
\begin{align}
 z_+(t,y)&=\frac{\mu\sin(2(y-ct))+c\cos(2(y-ct))}{4D},\label{eq:zplus}\\
 z_-(t,y)&=\frac{\mu\sin(2(y+ct))-c\cos(2(y+ct))}{4D}.\label{eq:zminus}
\end{align}
两支特征波的振幅均为
\begin{equation}\label{eq:amplitude}
 \boxed{\mathcal A=\frac1{4\sqrt{\mu^2+c^2}}>0}.
\end{equation}
其中，$\mathcal A$ 是 $z_+$ 或 $z_-$ 的空间振幅，即对 $y\in\T$ 的最大绝对值；它不同于压力常数 $A$，原解中相应声学分量的领先振幅为 $\eps^2\mathcal A$。
密度系数和速度系数的时间振幅分别为 $1/(4c\sqrt D)$ 和 $1/(4\sqrt D)$。不能用某一个时刻 $Q(t)=0$ 推断声学态消失。恒等式
\begin{equation}\label{eq:nonzero}
 c^2B_\infty(t)^2+Q_\infty(t)^2=\frac1{16D}
\end{equation}
给出整个声学相空间的统一非零下界。

二阶只依赖 $p'(1)=A\gamma$，不单独依赖 $p''(1)$。当 $\mu$ 很大时声学振幅约为 $1/(4\mu)$，其减弱来自源项很快被耗散；当 $\mu\downarrow0$ 时上述系数有有限极限，但源项消退时间 $1/\mu$ 发散，不能交换 $t\to\infty$ 与 $\mu\to0$。

\subsection{存在的抵消结构与不存在的抵消}
\begin{itemize}
\item $a=0$ 消掉了二阶密度--速度通量与二阶压力非线性；$\avg{U_1U_2}=0$，所以测试初值的二阶水平动量零模为零。
\item $W(x+\pi,y)=-W(x,y)$。由平移不变性和唯一性，参数 $-\eps$ 的解是参数 $\eps$ 的解水平平移 $\pi$，故水平平均量关于 $\eps$ 为偶函数。在具有足够参数可微性、允许损失 Sobolev 导数的固定时间展开中，平均量的三阶系数因此为零。本文保守的 $O(\eps^3)$ 误差并未用此改进。
\item 正负 $k_1$ 的相互作用生成 $k_1=0$，其 $k_2=\pm2$ 声学投影由 \eqref{eq:zplus}--\eqref{eq:amplitude} 确认非零。没有导致二次声学尾波完全消失的 null structure。
\end{itemize}

\section{精确平均方程：闭合误差不能删除}
只要密度为正，定义
\[
 r=\avg\rho,\qquad V=\frac{\avg{\rho u_2}}r,\qquad
 \eta=\rho-r,\quad w=u_2-V.
\]
其中，$r=r(t,y)$ 是水平平均密度，$V=V(t,y)$ 是质量加权的平均垂直速度，$\eta=\eta(t,x,y)$ 是相对于水平平均的密度振荡，$w=w(t,x,y)$ 是相对于 $V$ 的垂直速度振荡；这里的 $w$ 不是一维黎曼不变量 $w_\pm$。本节的 $r$ 对一般二维解定义，只在一维不变子类中与此前的一维密度完全相同。
注意 $\avg\eta=0$、$\avg{\rho w}=0$，但一般 $\avg w\ne0$。由守恒式严格得到
\begin{equation}\label{eq:exactmean}
 \boxed{\begin{aligned}
 r_t+(rV)_y&=0,\\
 (rV)_t+(rV^2+p(r))_y&=-\partial_y\mathcal S,\\
 \mathcal S&=\mathcal R+\Pi,\\
 \mathcal R&=\avg{\rho(u_2-V)^2},\qquad
 \Pi=\avg{p(\rho)}-p(r).
 \end{aligned}}
\end{equation}
其中，$\mathcal R$ 是垂直速度振荡的密度加权二次应力，$\Pi$ 是平均压力与平均密度所对应压力之差，$\mathcal S=\mathcal R+\Pi$ 是平均动量方程的总修正应力；它们都是 $t,y$ 的函数，$-\partial_y\mathcal S$ 是其对平均声学系统的作用力。
黏性平均为零是因为 $\avg{\partial_x^2u_2}=0$；这不表示其对 $\mathcal R,\Pi$ 的动力学无影响。由于 $p$ 凸，$\mathcal R,\Pi\ge0$，但 $\partial_y\mathcal S$ 没有符号。
等价的速度方程为
\[
 V_t+VV_y+\frac{p'(r)}r r_y=-\frac1r\mathcal S_y.
\]
精确 Taylor 形式
\[
 \Pi=\avg{\eta^2\int_0^1(1-\theta)p''(r+\theta\eta)\dd\theta}
\]
其中，$\theta\in[0,1]$ 是 Taylor 余项中的辅助积分参数，$p''$ 是压力对密度的二阶导数；$r+\theta\eta$ 是从平均密度 $r$ 到点值密度 $\rho$ 的线性插值。这里的 $\theta$ 与后文声学轮廓的空间相位变量仅共用字母。
这一恒等式说明压力修正由密度振荡决定，不能仅用 $p(r)$ 替代平均压力。

测试初值在固定时间展开中给出
\[
 \mathcal R=\frac{\eps^2}{2}e^{-2\mu t}\cos^2y+O(\eps^3),\qquad
 \Pi=O(\eps^4).
\]
其中，$O(\eps^3)$、$O(\eps^4)$ 表示在固定时间区间及所选足够高阶空间范数下，余项分别由与 $\eps$ 无关的常数乘 $\eps^3$、$\eps^4$ 控制；这些记号本身不保证增长时间区间上的一致估计。
利用偶性可在更高阶展开中改进前一式余项；这不影响二阶源。形式四阶压力主项为
\[
 \frac{\eps^4p''(1)H(t)^2}{4}.
\]
它在 $y$ 上恰为常数，故其在四阶平均动量中的导数为零；更高阶及一般初值不具有此简化。

如还需水平平均速度，定义 $J=\avg{\rho u_1}$、$W_m=J/r$，则
\[
 J_t+\partial_y(rW_mV+\mathcal C)=0,\qquad
 \mathcal C=\avg{\rho(u_1-W_m)(u_2-V)}.
\]
其中，$J=J(t,y)$ 是水平平均的水平动量密度，$W_m=W_m(t,y)$ 是质量加权的平均水平速度，$\mathcal C$ 是水平、垂直速度振荡的密度加权交叉应力；$W_m$ 的下标 $m$ 在这里表示平均量，不是后文 Sobolev 正则性阶数，$W_m$ 也不同于初始形状 $W$。
这显示慢子空间并不只含两支声学波。对本测试初值，水平反射 $x\mapsto-x$ 下 $u_1$ 为奇函数、$\rho,u_2$ 为偶函数，唯一性保持这一已有对称性，故 $J\equiv0$。这不是对模型另加的初值假设；推广至一般数据时必须处理 $J$。

\section{一般涡旋数据的二次声学选择映射}
这一节推广的是待研究的数据类，明确增加了“常密度、零水平平均、散度零”这三个条件；它们正是给定测试初值已有的性质。设 $u_0=\eps W$、$\avg W=0$、$\diver W=0$、$\rho_0=1$。则一阶仍为
$U(t)=e^{\mu t\partial_x^2}W$、$a=0$。写二阶平均密度和垂直速度为 $b_m,q_m$，则
\begin{equation}\label{eq:generalmean}
 (b_m)_t+(q_m)_y=0,\qquad
 (q_m)_t+c^2(b_m)_y=F(t,y),\quad
 F=-\partial_y\avg{U_2^2}.
\end{equation}
其中，$b_m=\avg b$、$q_m=\avg{Z_2}$ 分别是二阶水平平均密度、垂直速度系数，二者下标 $m$ 均表示平均；$F=F(t,y)$ 是一阶垂直速度 $U_2$ 的平方通过水平平均产生的声学源项。这里 $W$ 可取本节规定数据类中的一般初始形状，$e^{\mu t\partial_x^2}$ 是只沿水平变量演化的热半群，即方程 $\partial_tU=\mu\partial_x^2U$ 的解算子。
因此 $z_\pm=q_m\pm cb_m$ 满足 $(\partial_t\pm c\partial_y)z_\pm=F$，其散射轮廓是明确的二次映射
\begin{equation}\label{eq:scattermap}
 \boxed{\mathcal Q_\pm[W](\theta)
 =-\partial_\theta\int_0^\infty
 \avg{(e^{\mu s\partial_x^2}W_2)^2}(\theta\pm cs)\dd s.}
\end{equation}
其中，$\mathcal Q_\pm[W]$ 是由初始形状 $W$ 决定的两支二阶声学散射轮廓，方括号表示它对初始数据的依赖；$\theta\in\T$ 是沿相应声学特征移动的空间坐标，$s\ge0$ 是积分中的时间变量，$W_2$ 是 $W$ 的垂直分量，$\partial_\theta$ 是对轮廓空间坐标求导。$\mathcal Q_\pm$ 与单个测试模式的时间系数 $Q(t)$ 不是同一个量。
此式对展开系数严格成立。热半群在非零水平频率上指数衰减，故积分绝对收敛。Fourier 约定 $W_2=\sum_{k,\ell}\widehat W_{2;k,\ell}e^{i(kx+\ell y)}$ 下
\begin{equation}\label{eq:multiplier}
 \widehat{\mathcal Q_\pm[W]}(n)
 =-in\sum_{k\ne0}\sum_{\ell\in\mathbb Z}
 \frac{\widehat W_{2;k,\ell}\widehat W_{2;-k,n-\ell}}
 {2\mu k^2\mp icn}.
\end{equation}
其中，$k\in\mathbb Z\setminus\{0\}$ 和 $\ell\in\mathbb Z$ 是输入模式的水平、垂直整数波数，$n\in\mathbb Z$ 是输出声学轮廓的垂直波数；这里的 $k$ 是单个整数，区别于前面线性谱中的二维向量 $k=(k_1,k_2)$。$\widehat W_{2;k,\ell}$ 是 $W_2$ 的二维 Fourier 系数，$\widehat{\mathcal Q_\pm[W]}(n)$ 是输出轮廓的一维 Fourier 系数；所用系数约定分别以 $(2\pi)^{-2}$ 和 $(2\pi)^{-1}$ 归一化积分。分母中的 $\mp$ 与 $\mathcal Q_\pm$ 配对取相反符号。
这给出频率选择律、耗散--传播相位的竞争，也说明 $n=0$ 不被生成。对 $n\ne0$，乘子绝对值不超过 $1/c$；配合二维 Sobolev 代数估计，二次映射在足够高正则性中连续。测试数据已证明该映射不恒为零，因而其附近（在上述受约束数据类内）存在声学输出非零的开集。
不过瞬时 $F(0)\ne0$ 对一般数据不保证积分 \eqref{eq:scattermap} 非零，必须检查整个受迫历程的相位抵消。

\section{初步文献定位与逐项排除}
检索日期为 2026 年 10 月 30 日。使用的检索组合包括 compressible Navier--Stokes、anisotropic / horizontal / transverse viscosity、partial dissipation、shock formation、zero mode、mean flow、effective acoustics 和 Taylor--Green vortex。以下列出的定理号、模型或章节均核对了原始论文的相关正文；没有把摘要当作完整证明审查。2026 年论文按预印本处理，不把它们当作已经同行评审的结论。

\subsection{紧凑对照}
\small
\begin{longtable}{@{}>{\raggedright\arraybackslash}p{.20\textwidth}>{\raggedright\arraybackslash}p{.17\textwidth}>{\raggedright\arraybackslash}p{.30\textwidth}>{\raggedright\arraybackslash}p{.25\textwidth}@{}}
\toprule
论文与定位 & 定义域 & 黏性／压力 & 结论及与本题的差别\\
\midrule
\endfirsthead
\toprule
论文与定位 & 定义域 & 黏性／压力 & 结论及与本题的差别\\
\midrule
\endhead
CGM \cite{CGM}，Theorems 2--3 & $\R^2$ & 标量 $u_t+uu_x=u_{yy}$；无密度、压力 & 梯度爆破构造；不是可压缩系统或环面零模生成\\[4pt]
Bresch--Jabin \cite{BJ}，Theorem 3.2 & $\T^d$ & 小各向异性椭圆应力及体黏性；一般增长压力律 & 有限能量全球弱解；不覆盖本题退化算子\\[4pt]
Gao--Peng--Wu--Yao \cite{GPWY}，Theorems 1.1--1.2 & $\R^3_+$，滑移边界 & $\Delta_hu+\nu\partial_3^2u+\nabla\diver u$；$P=\rho^3/3$ & 一致正则性与 $\nu\to0$；极限仍有体黏性\\[4pt]
Feng--Hong--Wang--Wu \cite{FHWW}，Theorems 1.1--1.2 & $\R^3$ & $\mu\Delta_hu+(\mu+\zeta)\nabla\diver u$；$P'>0$ & 小解全球稳定及衰减；体黏性系数严格正\\[4pt]
Masmoudi--Sun--Wang--Zhang \cite{MSWZ}，Theorems 1.1--1.2 & $\T^2\times[0,a_3]$，Dirichlet & $\mu_1(\Delta_h+\nu\partial_z^2)u+\mu_2\nabla\diver u$；$\rho^\gamma/\gamma$ & 低 Mach、消失垂直黏性极限；固定时间近似与余项\\[4pt]
Huang--Li--Xu \cite{HLX}，Theorems 2.2--2.3 & $\T\times\R$ & 完整黏性及体黏性；$P'(1)=1$ & Couette 稳定、扩散波；零模也有扩散\\[4pt]
Antuono--Marrone \cite{AM}，\S\S2--4 & 二维双周期 & 速度方程 $\nu\Delta u$；线性化弱可压压力 & TG 涡声学修正、形式展开及数值对照\\[4pt]
Luk--Speck \cite{LS}，Theorem 16.1 & $\R\times\T$ & 无黏；真正非线性的正压律 & 平面 Euler 激波对有涡二维扰动的稳定性\\
\bottomrule
\end{longtable}
\normalsize

\subsection{与题目新颖性最有关的核对结果}
\paragraph{各向异性不等于方向完全退化。}
Bresch--Jabin 的 Theorem 3.2，正文 pp.~594--595、式 (3.6)--(3.9)，要求相对各向同性背景的小各向异性及耗散正性；数据是可含真空的有限能量数据，压力增长指数还有下界限制。它是弱解存在性结果，不能回答本题的正密度经典解激波。\cite{BJ}

Gao 等的式 (1.1)、极限式 (1.5) 明确保留 $\nabla\diver u$。Theorem 1.1 使用高阶 conormal 及负水平阶数等小能量数据，Theorem 1.2 给极限全球解；空间是半空间，边界为 $u_3=\partial_3u_h=0$。Feng 等的式 (1.1) 则要求 $\mu>0,3\zeta+2\mu>0$，所以 $\mu+\zeta>0$；Theorem 1.1 针对 $N\ge3$ 的小 $H^N$ 数据。两者的纵向声学耗散均不能移除后原样引用。\cite{GPWY,FHWW}

\paragraph{有效声学和平均／非平均模方法已有重要基础。}
Masmoudi 等取 Mach 数 $\epsilon\to0$、$\nu=\epsilon^\kappa$，Theorem 1.2 的范围含 $0<\kappa<3$、相容性与几何小除数条件；Appendix B.2 的 $(\mathrm{DB})_m^\alpha$ 是耦合阻尼 Burgers 声学方程，Lemma B.3 给局部适定性。其近似解加余项稳定性的组织可借鉴，但参数和时间尺度均不同于本题固定声速的 $\eps^{-2}$ 演化。\cite{MSWZ}

Huang--Li--Xu 的式 (2.4)--(2.5) 使用平均守恒变量和非零模应力，式 (2.19)、(2.24) 构造 Burgers 扩散波及耦合修正；Theorems 2.2--2.3 针对 Couette 背景下 $O(\mu^\alpha)$、$\alpha>11/3$ 的扰动及附加可积性／反导数条件。这里平均声学模有纵向扩散，不能据此得到本题无阻尼零模的长期行为。\cite{HLX}

\paragraph{Taylor--Green 涡产生声波不是新的物理观察。}
Antuono--Marrone 的 \S2 式 (1) 采用速度方程中常运动黏度 $\nu\Delta u$；弱可压缩参数为 $\mathrm{Ma}^2$，压力近似线性，\S3 构造 Taylor--Green 修正并与数值结果比较。原文讨论了不可压缩 TG 初值的声学污染。\cite{AM}
本题潜在新意应放在\emph{耗散核内的非零残留声学数据及其定量长期演化}。访问限制须说明：该文出版社及 CNR PDF 直取遇到 403，核对依赖检索工具返回的原始 PDF 正文片段（\S2 及 \S3 的相关段落），未审阅全文证明；不为它编造严格误差定理或定理号。

\subsection{CGM 的工具和可压缩系统的实质障碍}
CGM 原文 v2 的 Theorems 2--3（PDF pp.~4--5）构造局域化光滑初值的梯度奇性。式 (1.6)--(1.7) 利用对称轴得到
$\xi=-u_x|_{x=0}$ 满足 $\xi_t=\xi^2+\xi_{yy}$。\S1.5(2) 讨论对称性破缺与倾斜；\S5.2 的稳定分析使用奇偶结构，不能据定理中的简写把它读成任意二维扰动稳定性。\cite{CGM}

下面是对本题的独立结构比较。可迁移的思路包括：沿奇性方向的展开、横向抛物调制、自相似变量和加权能量、内外区匹配。难点在于本系统有密度和两个声学族，梯度方程含压力 Hessian，速度黏性是 $\mu/\rho$；没有上述标量轴上闭合。

倾斜的危险可直接计算：对标量方程取 $u(t,x,y)=F(t,x-\delta y)$，便得到
\[
 F_t+FF_z=\delta^2F_{zz}.
\]
这说明横向黏性可能因波面方向改变而作用于压缩方向。对本题要另行研究这种几何效应，而不能把它视为 Euler 激波方程中的一般低阶误差。

\subsection{Euler 激波参照和线性耗散判据}
Luk--Speck 的 Theorem 16.1，式 (16.0.66)、(16.0.69)，处理满足相对小量条件的简单平面波扰动，允许非零涡量；声学几何聚焦使密度和速度梯度失界而状态量保持有界。它提供处理多速率和涡量的框架，但没有本题的方向退化二阶算子，也没有由纯涡旋经过长时间产生弱压缩的过程。\cite{LS}

Shizuta--Kawashima 的 Theorem 1.1 将对称线性双曲--抛物系统的严格耗散与代数耦合条件联系起来，Theorem 1.2 给相应衰减。\cite{SK}
本题在方向 $k=(0,k_2)$ 上扩散符号为零且有声学特征向量，显然不满足全方向严格耗散条件。这一独立检验解释了为何不能套用以该条件为前提的全球小解定理。

\subsection{查重边界}
以上文献没有直接给出本题的零模交接、慢时间调制或真正二维激波结论。\textbf{这只排除了这些具体定理的直接适用性，不证明题目尚未解决。}“partial viscosity”可能仅指密度变量没有扩散；“degenerate viscosity”也常指真空处密度依赖黏度退化，都不等于本题的空间方向退化。
下一轮须追踪 TG 声学修正、非线性几何光学、退化双曲--抛物中心模和横向黏性激波的引用链，尤其核查是否已有抽象定理直接包含本文的二次选择映射及慢时间结论。没有完成这一步前，不宜在论文摘要中写“首次”或“解决公开问题”。

\section{第一步可以做到什么：定量展开与对数时间交接}
以下给出一个可由局部能量法闭合的命题及证明。它只研究给定测试初值；不把有限时间的 Taylor 公式延伸到未知的激波时间。这里的推导是数学证明方案，局部构造和交换子估计写到可复核的能量层面；它不是已经发表或已经同行评审的定理。

\begin{proposition}[二阶声学交接与寿命下界]\label{prop:matching}
固定 $A,\gamma,\mu>0$（$\gamma>1$）、整数 $m\ge6$ 和 $L>1/\mu$，令
\[
 t_\eps=L\log(1/\eps).
\]
其中，$m$ 是用于能量估计的空间导数阶数，$L$ 是与 $\eps$ 无关且大于 $1/\mu$ 的固定常数，$t_\eps$ 是水平振荡衰减后与声学近似进行比较的交接时间，$\log$ 是自然对数。
对充分小的 $\eps>0$，给定测试初值的光滑解至少存在到 $t_\eps$，且 $1/2<\rho<3/2$。
令 $(b,Z)$ 由 \eqref{eq:fullsecond} 定义。在 $0\le t\le t_\eps$ 上有
\begin{equation}\label{eq:matchingerror}
 \norm{\rho-1-\eps^2b}_{H^m}
 +\norm{u-\eps e^{-\mu t}W-\eps^2Z}_{H^m}
 \le C\bigl(\eps^3+\eps^4t\bigr)e^{C\eps t}.
\end{equation}
其中，$H^m=H^m(\T^2)$ 是函数及其直到 $m$ 阶的所有弱空间导数均平方可积的 Sobolev 空间，$\norm{\cdot}_{H^m}$ 是相应范数；向量或多分量函数的范数按各分量范数的平方和开方计算。$C$ 是可随估计改变的正数，可依赖固定的 $A,\gamma,\mu,m,L$ 及测试轮廓，但不依赖 $\eps$ 和所述区间内的 $t$。下文 $O(\,\cdot\,)$ 表示所讨论范数至多为括号内量的固定常数倍，$o(\,\cdot\,)$ 表示该范数与括号内量之比在 $\eps\downarrow0$ 时趋于零。
进而
\begin{equation}\label{eq:handoff}
 \eps^{-2}(\rho-1,u)(t_\eps)
 -(B_\infty(t_\eps)\cos2y,\ 0,\ Q_\infty(t_\eps)\sin2y)
 \longrightarrow0 \quad\text{于 }H^m(\T^2).
\end{equation}
其中，$(\rho-1,u)$ 按 $(\rho-1,u_1,u_2)$ 的顺序组成三分量向量，右侧比较向量的中间分量 $0$ 表示没有领先水平速度；箭头表示当 $\eps\downarrow0$ 时，整个差值的 $H^m$ 范数趋于零。
特别地，此时完整解的扰动大小为 $O(\eps^2)$，联合声学分量有非零的 $O(\eps^2)$ 下界。由重启局部能量估计，还得到某个 $c_*>0$ 使经典寿命
\begin{equation}\label{eq:lowerlife}
 T_{\max}(\eps)\ge t_\eps+c_*\eps^{-2}.
\end{equation}
其中，$T_{\max}(\eps)$ 是振幅为 $\eps$ 的给定初值所对应经典解的最大存在时间，允许其为 $+\infty$；$c_*$ 是依赖固定参数和估计阶数、但不依赖 $\eps$ 的正寿命常数，与背景声速 $c$ 不同。
这只是寿命下界，不是寿命渐近或有限时间爆破上界。
\end{proposition}

\subsection{对称能量与局部构造的关键估计}
使用声速变量
\[
 q=\frac{2(c(\rho)-c)}{\gamma-1},\qquad
 c(\rho)=c+\frac{\gamma-1}{2}q.
\]
其中，$q$ 是相对于背景声速 $c$ 归一化的声速扰动，用来对称化双曲主部；它不是前文的二阶平均垂直速度系数 $q_m$。
在正密度区间内，这是光滑可逆变换，且
\begin{equation}\label{eq:symmetric}
 \begin{aligned}
 q_t+u\cdot\nabla q+c(\rho)\diver u&=0,\\
 u_t+u\cdot\nabla u+c(\rho)\nabla q&=\frac\mu{\rho(q)}\partial_x^2u.
 \end{aligned}
\end{equation}
其中，$\rho(q)$ 是由上述声速变量变换反解得到的密度，即 $\rho(q)=[1+(\gamma-1)q/(2c)]^{2/(\gamma-1)}$，$\mu/\rho(q)$ 是速度方程中的变系数水平扩散系数。
双曲主部对称，扩散矩阵半正定。设 $E_m\simeq\norm{q,u}_{H^m}^2$，对充分小解可取与标准范数等价的能量，使
\begin{equation}\label{eq:energyhigh}
 E_m'+c_\mu\norm{\partial_xu}_{H^m}^2
 \le C(E_m^{1/2}+E_m)E_m.
\end{equation}
其中，$E_m=E_m(t)$ 是控制 $(q,u)$ 的 $m$ 阶空间导数的能量，$E_m'$ 是其时间导数；$\simeq$ 表示存在两个与 $\eps,t$ 无关的正数给出双边比较，因而 $E_m$ 与 $\norm{q,u}_{H^m}^2$ 等价。$c_\mu>0$ 是吸收扩散交换子后保留的耗散系数，可依赖 $\mu$ 及密度的固定正上下界，$\norm{\partial_xu}_{H^m}^2$ 是受控的水平耗散量。
核对该式的主要步骤如下。对称性消去线性声学交叉项；输运和可变声速项由 Sobolev 交换子估计控制。对 $a(q)=\mu/\rho(q)$，扩散交换子用
\[
 \norm{[D^\alpha,a]\partial_x^2u}_{L^2}
 \le C\bigl(\norm{\nabla a}_{L^\infty}\norm{\partial_xu}_{H^m}
       +\norm{a-\mu}_{H^m}\norm{\partial_x^2u}_{L^\infty}\bigr),
 \quad |\alpha|\le m,
\]
其中，$\alpha=(\alpha_1,\alpha_2)$ 是非负整数多重指标，$|\alpha|=\alpha_1+\alpha_2$ 是总阶数，$D^\alpha=\partial_x^{\alpha_1}\partial_y^{\alpha_2}$ 是相应空间导数，$[D^\alpha,a]f=D^\alpha(af)-aD^\alpha f$ 是微分与乘法的交换子，$f$ 是任意足够光滑的待作用函数。这里 $a=a(q)$ 是扩散系数函数，不是前文密度展开的一阶系数 $a$；$L^2$ 是平方可积函数空间，$L^\infty$ 是本质有界函数空间，相关范数均取在 $\T^2$ 上。
第一项用水平耗散吸收；第二项利用
\[
 H^m(\T^2)\hookrightarrow W^{2,\infty}(\T^2).
\]
其中，$W^{2,\infty}$ 是函数及其直到二阶的弱空间导数均本质有界的 Sobolev 空间，$\hookrightarrow$ 表示连续嵌入，即左侧范数控制右侧范数。主扩散项积分分部产生的 $a_x=\partial_xa(q)$ 项也由耗散及 $\norm{a_x}_{L^\infty}^2E_m$ 控制。不会由此要求 $\partial_yD^mu$ 的耗散；这里 $D^mu$ 泛指 $u$ 的任意 $m$ 阶空间导数。更高阶光滑性通过 tame 高阶估计延拓，其增长系数由已控制的低阶 $H^m$ 范数给出；不对每个高阶范数另取越来越小的 $\eps$ 阈值。

对正密度的小初值，频率截断或暂加正则化得到近似解，\eqref{eq:energyhigh} 给参数一致局部估计，紧性和差值估计给存在唯一性与延拓。它在初值范数 $O(\eps)$ 时给出
\[
 T_{\max}(\eps)\ge c_0/\eps,\qquad
 \norm{q,u}_{H^m}\le C\eps\quad(0\le t\le c_0/\eps).
\]
其中，$c_0>0$ 是局部能量估计给出的、与 $\eps$ 无关的存在时间常数，与声速 $c$ 及重启后的常数 $c_*$ 分别记号。这已覆盖任意固定 $L$ 的对数区间。构造中的辅助正则化不属于最终模型。

\subsection{残差计算与证明闭合}
设
\[
 \rho^a=1+\eps^2b,\qquad u^a=\eps U+\eps^2Z.
\]
其中，$\rho^a,u^a$ 是截断到二阶的近似密度和近似速度，上标 $a$ 是“近似”的标记，不表示幂次，也不表示上述扩散系数函数。
在原始质量方程中的残差精确为
\[
 R_\rho=\eps^3\diver(bU)+\eps^4\diver(bZ).
\]
其中，$R_\rho$ 是将 $(\rho^a,u^a)$ 代入质量方程后得到的残差，即 $\partial_t\rho^a+\diver(\rho^au^a)$。
在速度方程中的三阶残差为
\[
 \eps^3\{U\cdot\nabla Z+Z\cdot\nabla U+\mu b\partial_x^2U\}.
\]
其余为 $O(\eps^4)$，包括压力展开、$Z\cdot\nabla Z$ 和变系数黏性修正。因为 $b,Z$ 及其任意固定阶空间导数有界，而 $U$ 指数衰减，
\begin{equation}\label{eq:residual}
 \norm{R_\rho,R_u}_{H^m}\le C(\eps^3e^{-\mu t}+\eps^4).
\end{equation}
其中，$R_u$ 是将 $(\rho^a,u^a)$ 代入速度方程 \eqref{eq:primitive} 后用左端减去右端得到的向量残差，$\norm{R_\rho,R_u}_{H^m}$ 同时度量质量方程和速度方程的近似误差。
将近似解也变换为声速变量，并定义差值
\[
 q^a=\frac{2(c(\rho^a)-c)}{\gamma-1},\qquad
 e=(q-q^a,u-u^a).
\]
其中，$q^a$ 是近似密度对应的声速扰动，$e$ 是精确解与近似解在对称变量中的差值向量，不是自然指数函数的底数。对 $e$ 的对称差值能量使用上一小节同样的扩散交换子估计，给出
\[
 \frac{\mathrm d}{\mathrm dt}\norm{e}_{H^m}
 \le C\eps\norm{e}_{H^m}+C(\eps^3e^{-\mu t}+\eps^4),\qquad e(0)=0
\]
其中，$e(0)=0$ 表示近似解与精确解具有相同初值，指数项 $e^{-\mu t}$ 表示一阶速度的衰减因子。
零范数处可先正则化能量的平方根。这里近似解所需的更高导数均由显式有限 Fourier 展开控制；如果用一般局部稳定性引理，可以在高两阶范数构造解，再在 $H^m$ 比较。
积分即得 \eqref{eq:matchingerror}，再由变量变换局部等价性回到 $(\rho,u)$。

在 $t=t_\eps$，$\eps e^{-\mu t_\eps}=\eps^{1+\mu L}=o(\eps^2)$；$H,P$ 以及平均系数的指数瞬态都趋零；\eqref{eq:matchingerror} 除以 $\eps^2$ 也趋零。因此得到 \eqref{eq:handoff}。由 \eqref{eq:nonzero}，这个声学态不会沿某个相位子列整体消失。由于 $t_\eps$ 的声学相位随 $\eps$ 转动，这不是收敛到一个固定空间函数，而是与给定时刻的显式自由轨道接近。

最后在 $t_\eps$ 重启 \eqref{eq:energyhigh}，此时 $E_m^{1/2}\le C\eps^2$，就得到 \eqref{eq:lowerlife}。相同论证在较高阶施行，故这一重启不是只控制低阶范数的形式操作。\hfill$\square$

固定时间结论是该命题的弱形式：对每个固定 $T$，有
\[
 \sup_{0\le t\le T}\norm{e(t)}_{H^m}=O_T(\eps^3).
\]
其中，$T$ 是与 $\eps$ 无关的固定时间上界，$\sup$ 表示对该时间区间取上确界，$O_T$ 的下标表示估计常数允许依赖 $T$。仅有这条弱形式时，不能令 $T=T(\eps)\to\infty$ 而忽略常数的增长；上面的显式时间估计正是补上的环节。

\section{最值得继续证明的论文命题：激波前的慢时间声学调制}
命题 \ref{prop:matching} 越过了瞬时生成，但本质上仍是精细局部扰动理论。更有希望成为严格论文主结果的是下述明确命题：\emph{从无初始声学分量的二维涡旋出发，在 $O(\eps^{-2})$ 的非线性时间上，严格识别耗散核内的有效声学演化。}

\subsection{形式调制系统及其系数}
在一维平均 Euler 系统中使用中心化黎曼变量
\[
 j_\pm=V\pm\frac{2(c(r)-c)}{\gamma-1}.
\]
其中，$j_+,j_-$ 是平均密度 $r$ 和质量加权垂直速度 $V$ 对应的两支中心化黎曼变量，在平衡态 $(r,V)=(1,0)$ 时均为零；正、负号分别对应背景传播速度为 $+c$ 和 $-c$ 的声学族。
保留真实平均应力时，它们满足精确方程
\begin{equation}\label{eq:riemannforced}
 \left[\partial_t+\left(\pm c+\kappa j_\pm+\beta j_\mp\right)\partial_y\right]j_\pm
 =-\frac{\mathcal S_y}{r},\quad
 \kappa=\frac{\gamma+1}{4},\quad\beta=\frac{3-\gamma}{4}.
\end{equation}
其中，$\kappa$ 是同一声学族的自相互作用系数，$\beta$ 是相反声学族之间的耦合系数，$j_\mp$ 指与 $j_\pm$ 相反的一族，$\mathcal S_y=\partial_y\mathcal S$ 是精确平均应力的垂直导数。
取慢时间与随声学特征移动的坐标
\[
 \tau=\eps^2t,\qquad \theta=y\mp ct,\qquad
 j_\pm\sim\eps^2f_\pm(\tau,y\mp ct).
\]
其中，$\tau$ 是描述 $O(\eps^{-2})$ 物理时间演化的慢时间，$\theta$ 是沿相应声学族移动的周期空间坐标，$f_\pm(\tau,\theta)$ 是归一化到 $O(1)$ 大小的声学包络；此处 $\sim$ 表示形式领先阶近似，尚不表示已经证明的误差估计。
两支初始声学轮廓由二次选择映射给出：
\[
 f_+(0,\theta)=\frac{\mu\sin2\theta+c\cos2\theta}{4D},\qquad
 f_-(0,\theta)=\frac{\mu\sin2\theta-c\cos2\theta}{4D}.
\]
在互逆声学族的非共振项被正常形或平均法严格控制之后，\emph{预期}领先包络满足
\begin{equation}\label{eq:burgersmod}
 (f_\pm)_\tau+\kappa f_\pm(f_\pm)_\theta=0.
\end{equation}
这里没有黏性项，因为领先波完全位于 $k_1=0$。反向波具有零均值，所以其沿另一族特征的平均为零；这一观察解释 \eqref{eq:burgersmod}，但还不是对原二维系统的平均化定理。
在一维 Euler 内部，当 $\gamma=3$ 时 $\beta=0$，两支黎曼方程恰好解耦。把 $\gamma$ 限制为 $3$ 会是额外限制，可作为后续技术试验，不能冒充本轮所有 $\gamma>1$ 的结论。

\subsection{候选主定理的具体表述（待证明）}
令 $f_\pm$ 解 \eqref{eq:burgersmod}，并定义
\[
 b_{\rm eff}(t,y)=\frac{f_+(\eps^2t,y-ct)-f_-(\eps^2t,y+ct)}{2c},\qquad
 v_{\rm eff}(t,y)=\frac{f_+(\eps^2t,y-ct)+f_-(\eps^2t,y+ct)}2.
\]
其中，$b_{\rm eff}$ 是由两支包络重建的有效密度扰动系数，$v_{\rm eff}$ 是有效垂直速度系数，下标 $\mathrm{eff}$ 表示“有效”；对应的物理扰动分别是 $\eps^2b_{\rm eff}$ 和 $\eps^2v_{\rm eff}$。
对于每个固定 $0<\tau_0<\tau_*$，希望证明原二维解存在到 $\tau_0\eps^{-2}$，密度保持在一个固定正区间内，并且对一个明确的整数 $m$，
\begin{equation}\label{eq:target}
 \sup_{t_\eps\le t\le\tau_0\eps^{-2}}
 \eps^{-2}\norm{(\rho-1,u)-\eps^2(b_{\rm eff},0,v_{\rm eff})}_{H^m(\T^2)}
 \longrightarrow0.
\end{equation}
其中，$\tau_0$ 是严格小于有效 Burgers 首次激波慢时间 $\tau_*$ 的固定慢时间上界，$\tau_*$ 的具体值见 \eqref{eq:formalshock}；$\tau_0\eps^{-2}$ 是对应的物理时间上界，$\sup$ 要求误差在整个所列时间区间上一致趋于零，极限仍取 $\eps\downarrow0$。
应尽量将 $o(1)$ 改成显式幂次或对数误差率，并说明正则性损失和 $\tau_0\uparrow\tau_*$ 时常数的退化。定理首先针对给定测试数据；再推广到二次映射非退化的数据开集，需保留一般会产生的水平平均剪切。
这比“在一个短固定区间内有二阶项”更有论文价值，也比直接承诺二维激波稳妥。\eqref{eq:target} 即使成立，仍仅覆盖有效激波前的任意固定慢时间，不能自动推出原方程有限时间奇性。

\subsection{需要完成的关键估计}
\begin{enumerate}
\item \textbf{慢模与快模分解。}定义
\[
 P_0f=\avg f,\qquad P_\perp f=(I-P_0)f=f-\avg f.
\]
其中，$f$ 是待分解的标量或向量函数，$P_0$ 是取水平平均的投影，将所得平均视为不依赖 $x$ 的二维函数；$I$ 是恒等算子，$P_\perp$ 是取非零水平 Fourier 模的投影，作用于向量时逐分量进行。应在围绕时变一维声学背景的方程中控制 $P_\perp(\rho,u)$。需结合水平耗散与密度--速度补偿估计，证明在整个激波前慢时间内其应力对慢模的累计作用为 $o(\eps^2)$，并处理所需空间导数。
\item \textbf{平均应力的定量控制。}用 \eqref{eq:exactmean} 保留 $\mathcal R,\Pi$，例如在交接后建立足够强的
\[
 \int_{t_\eps}^{\tau_0\eps^{-2}}\norm{\mathcal S_y(t)}_{H^m}\dd t=o(\eps^2).
\]
其中，该积分度量交接后直到固定激波前慢时间 $\tau_0$ 的平均应力源的累计大小；这里 $\mathcal S_y$ 作为不依赖 $x$ 的二维函数取 $H^m$ 范数。也可给出等价的带正常形修正的估计。仅有每时刻 $O(\eps^4)$ 还不够：累计 $\eps^{-2}$ 时间可达 $O(\eps^2)$。
\item \textbf{两声学族的正常形。}未平均的交叉项 $\beta j_\mp\partial_yj_\pm$ 也是 $O(\eps^4)$。利用特征速度差 $2c$ 和零均值消掉非共振部分，保留自相互作用 Burgers 项，并控制导数损失。
\item \textbf{高阶近似与长期稳定。}单纯二阶线性近似的四阶残差积分会达到主项大小；需加入非线性一维演化或多尺度修正，使剩余源项在慢时间上足够小。系数随 $\tau_0$ 的控制必须写清楚。
\end{enumerate}

\section{激波时间猜想：合理但尚有重大缺口}
对 \eqref{eq:burgersmod}，两个初始轮廓均有
\[
 \min f_\pm'(0,\theta)=-\frac1{2\sqrt D}.
\]
其中，$f_\pm'$ 表示对周期空间坐标 $\theta$ 求导，$\min$ 是对 $\theta\in\T$ 取最小值，其负值刻画初始包络的最强压缩斜率。
所以有效 Burgers 的首次梯度爆破慢时间为
\begin{equation}\label{eq:formalshock}
 \tau_* =\frac{8\sqrt{\mu^2+A\gamma}}{\gamma+1}.
\end{equation}
其中，$\tau_*$ 是有效无黏 Burgers 包络首次出现梯度无界的慢时间，不是已经证明的原二维系统激波时刻。
这给出更具体、可检验的形式猜想
\[
 T_\eps\ \text{可能满足}\quad
 \eps^2T_\eps\longrightarrow\frac{8\sqrt{\mu^2+A\gamma}}{\gamma+1}.
\]
其中，$T_\eps$ 是所猜想的原二维解首次激波形成的物理时间；若这种激波确为首次经典奇性，则它等于 $T_{\max}(\eps)$。箭头表示猜想中的 $\eps\downarrow0$ 寿命渐近关系。
目前既未证明上界，也未证明这个常数。$O(\eps^2)$ 幅度推出 $O(\eps^{-2})$ 自陡化时间，依赖固定波长、真正非线性系数非零、存在压缩以及长期无有效阻尼等条件，不是仅从振幅即可推出的定律。

最主要的失败机制如下。
\begin{itemize}
\item \textbf{累积反馈。}高阶 Reynolds 应力、压力修正、生成的剪切模及两声学族相互作用可能在该时间尺度改变波形；必须估计，不能用“比主项高阶”排除。
\item \textbf{几何黏性。}若声学相位出现水平依赖，$\mu\partial_x^2$ 在相位坐标中含有二阶项
\[
 \mu(\partial_x\phi)^2\partial_\phi^2.
\]
其中，$\phi=\phi(t,x,y)$ 是描述声学波面位置的相位函数，$\partial_x\phi$ 度量波面对水平方向的依赖，$\partial_\phi$ 是保持其余新坐标不变时沿相位坐标求导的算子。即使系数很小，它在梯度急剧增大时也可能成为主导。CGM 的倾斜例说明这是一项结构性困难，而不是普通低阶扰动。环面周期性需另行检查，因此它不是给定初值的现成反例。
\item \textbf{临近奇性时的误差失控。}激波前固定 $\tau_0$ 的 Sobolev 逼近不能一直取到 $\tau_*$。要证明爆破，须建立压缩梯度的 Riccati 不等式或声学逆特征密度方程，并对水平二阶项作奇性适配估计。
\item \textbf{条件变化。}若让 $\mu$ 随 $\eps$ 变化，快慢分离可能消失；若换初值，一般二次映射可能为零，或产生不同波长和相位。作为直接反例，下式给出精确剪切解：
\[
 \rho=1,\qquad u=(0,\eps e^{\mu t\partial_x^2}f(x))
\]
其中，$f(x)$ 是任意光滑的 $2\pi$ 周期剪切轮廓，与前述声学包络 $f_\pm$ 不同；$e^{\mu t\partial_x^2}$ 是水平热半群，将初始轮廓演化为水平热方程的解。该解不生成声学波，但这不否定给定二维测试数据。
\end{itemize}

\section{下一步计算与投入建议}
建议继续，但把论文主目标暂定为 \eqref{eq:target} 的\textbf{激波前有效动力学}，把真正二维激波形成列为后续目标。只写一维嵌入或瞬时生成，不足以形成有竞争力的论文；一般“涡旋产生声波”的叙述也已有先例。

下一步按优先级完成：
\begin{enumerate}
\item 把命题 \ref{prop:matching} 的局部构造、差值能量与参数依赖写成完整引理，确认所有 Sobolev 指标及平均量的偶性。
\item 求三阶非零水平模，以及四阶平均方程中的全部项；区分可积的快模源与声学自相互作用，提取与 \eqref{eq:burgersmod} 一致的共振项。
\item 证明以一维非线性声学解为背景的非零水平模衰减和应力积分估计。先在固定 $\tau_0<\tau_*$ 上闭合，再研究常数接近奇性时如何退化。
\item 辅助数值若开展，应保持守恒动量与唯一的 $\mu\partial_x^2u$；记录声学特征轮廓、$\mathcal R,\Pi$、密度上下界、$\norm{\partial_yu_2}_\infty$ 及非零水平模能量，至少做空间分辨率和时间步长收敛。其中，范数下标 $\infty$ 是 $L^\infty(\T^2)$ 的简写，该量度量垂直速度的垂直梯度最大幅值。高频滤波或垂直人工黏性会直接改变问题，必须单独量化。
\end{enumerate}
本轮附带的计算脚本只验证二阶 ODE、非抵消恒等式与符号系数，不是原 PDE 激波的数值证据。13 项符号恒等式均通过；独立四阶 Runge--Kutta 积分在 $(\mu,c)=(1,1),(0.2,1),(3,1),(1,2)$、步长 $10^{-3}$、区间 $[0,8]$ 上与闭式解的最大绝对差小于 $7.5\times10^{-12}$。结果保存在 \texttt{verification.json}，不据此作激波数值判断。
如果后续发现倾斜黏性阻止稳定激波，仍可转向“耗散核内非线性声学选择与长时间调制”，或明确限定到 $\gamma=3$ 的压力律以隔离双声学族耦合；这些都是有数学内容的相邻问题。现在没有充分理由因为长期爆破困难而否定整个方向。

\clearpage
\begin{thebibliography}{99}
\raggedright
\bibitem{CGM}
C. Collot, T.-E. Ghoul, N. Masmoudi.
\emph{Singularity formation for Burgers equation with transverse viscosity}.
Annales scientifiques de l'\'Ecole normale sup\'erieure 55 (2022), 1047--1133.
\href{https://arxiv.org/pdf/1803.07826}{原始论文 arXiv:1803.07826v2}；
\href{https://doi.org/10.24033/asens.2513}{DOI:10.24033/asens.2513}.
核对位置：Theorems 2--3，式 (1.6)--(1.7)，\S1.5(2)，\S5.2。

\bibitem{BJ}
D. Bresch, P.-E. Jabin.
\emph{Global existence of weak solutions for compressible Navier--Stokes equations: Thermodynamically unstable pressure and anisotropic viscous stress tensor}.
Annals of Mathematics 188 (2018), 577--684.
\href{https://annals.math.princeton.edu/wp-content/uploads/annals-v188-n2-p04-p.pdf}{出版社原始 PDF}.
核对位置：Theorem 3.2，式 (3.6)--(3.9)，pp.~594--595。

\bibitem{GPWY}
Jincheng Gao, Lianyun Peng, Jiahong Wu, Zheng-an Yao.
\emph{Global uniform regularity and vanishing vertical viscosity limit for the compressible Navier--Stokes equations in the half-space}.
预印本 (2026).
\href{https://arxiv.org/html/2605.23233v1}{原始论文 arXiv:2605.23233v1}.
核对位置：式 (1.1)、(1.5)，Theorems 1.1--1.2，Remark 1.5。

\bibitem{FHWW}
Zefu Feng, Guangyi Hong, Yinghui Wang, Jiahong Wu.
\emph{Global stability and anisotropic large-time behavior of the three-dimensional compressible Navier--Stokes equations with eddy diffusion}.
预印本 (2026).
\href{https://arxiv.org/html/2607.06264v1}{原始论文 arXiv:2607.06264v1}.
核对位置：式 (1.1)，Theorems 1.1--1.2。

\bibitem{MSWZ}
N. Masmoudi, C. Sun, C. Wang, Z. Zhang.
\emph{Incompressible and vanishing vertical viscosity limit for the compressible Navier--Stokes system with Dirichlet boundary conditions}.
预印本 (2025).
\href{https://arxiv.org/html/2501.05063v1}{原始论文 arXiv:2501.05063v1}.
核对位置：Theorems 1.1--1.2，Appendix B.2，Lemma B.3。

\bibitem{HLX}
Feimin Huang, Rui Li, Lingda Xu.
\emph{Nonlinear stability threshold for compressible Couette flow}.
预印本 (2024).
\href{https://arxiv.org/html/2409.01142v1}{原始论文 arXiv:2409.01142v1}.
核对位置：Theorems 2.2--2.3，式 (2.4)--(2.5)、(2.19)、(2.24)。

\bibitem{AM}
M. Antuono, S. Marrone.
\emph{Weakly Compressible Approximation of the Taylor--Green Vortex Solution}.
Studies in Applied Mathematics 154 (2025), e12792 (online 2024).
\href{https://doi.org/10.1111/sapm.12792}{DOI:10.1111/sapm.12792}.
核对位置：\S2 式 (1) 及相关展开，\S3 相关正文片段；全文访问限制见正文。

\bibitem{LS}
J. Luk, J. Speck.
\emph{Shock formation in solutions to the 2D compressible Euler equations in the presence of non-zero vorticity}.
Inventiones mathematicae 214 (2018), 1--169.
\href{https://arxiv.org/pdf/1610.00737}{原始论文 arXiv:1610.00737}；
\href{https://doi.org/10.1007/s00222-018-0799-8}{DOI:10.1007/s00222-018-0799-8}.
核对位置：Theorem 16.1，式 (16.0.66)、(16.0.69)。

\bibitem{SK}
Y. Shizuta, S. Kawashima.
\emph{Systems of equations of hyperbolic-parabolic type with applications to the discrete Boltzmann equation}.
Hokkaido Mathematical Journal 14 (1985), 249--275.
\href{https://www.math.sci.hokudai.ac.jp/hmj/page/14-2/pdf/HMJ_14_2_1985_249-275.pdf}{出版社原始 PDF}.
核对位置：Theorems 1.1--1.2，pp.~250--251。
\end{thebibliography}

\end{document}
